\section{Vector Space of Linear Maps}

\subsection{Definition and Examples of Linear Maps}

% 3.1
\begin{mydef} [linear map]
  A \lm from $V$ to $W$ is a function $T:V\to W$ with the following properties:
  \begin{itemize}
    \item \emph{additivity:} $T(u+v)=Tu + Tv \qquad \forall u,v \in V$
    \item \emph{homogeneity:} $T(\lambda v)=\lambda (Tv) \qquad
    \forall \lambda \in \mathbb{F}, \; \forall v\in V$
  \end{itemize}
\end{mydef}

% 3.2
\begin{mydef} [notation $\linmap (V,W), \, \linmap (V)$]
  The set of all linear maps from $V$ to $W$ is denoted by $\mathcal{L}(V,W)$. For the set of all linear maps from $V$ to itself, we use $\mathcal{L}(V) :\equiv \mathcal{L}(V,V)$.
\end{mydef}

% 3.4
\setcounter{thm}{3}
\begin{thm}[linear map lemma]
  \label{thm: linear map lemma}
  Suppose $v_1, \dots, v_n$ is a basis of $V$ and $w_1, \dots, w_n \in W$, where the $w$'s do not have any specific properties. Then there exists a unique linear map $T \in \linmap(V,W)$ such that
  \begin{equation}
    Tv_{k} = w_k \quad \forall k \in \{1, \dots, n\}.
  \end{equation}

  It is of the form
  \begin{equation}
    T(c_1 v_1 + \cdots + c_n v_n) = c_1 w_1 + \cdots + c_n w_n.
  \end{equation}
\end{thm}
\begin{prf}
  Since the list $v_1, \dots, v_n$ is a basis of $V$, every $v \in V$ has a unique expansion
  \begin{equation}
    v = c_1 v_1 + \cdots + c_n v_n,
  \end{equation}

  where \( c_1, \dots, c_n \in \myF \). Define $T: V \to W$ by
  \[
    T(c_1 v_1 + \cdots + c_n v_n) := c_1 w_1 + \cdots + c_n w_n.
  \]
  By uniqueness of the basis expansion, this rule defines a well-defined function \(T\colon V\to W\) (otherwise a single \(v\) could receive different images).
  In particular, choosing \(v=v_k\) (so \(c_k=1\), \(c_i=0\) for \(i\neq k\)) gives
  \[
    T(v_k)=w_k.
  \]

  Next, we verify that \( T \) is linear. Let \( u = a_1 v_1 + \cdots + a_n v_n \in V \) and \( v = b_1 v_1 + \cdots + b_n v_n \in V \), and let \( \lambda \in \myF \). Then,
  \begin{align}
    T(u + v) &= T((a_1 + b_1)v_1 + \cdots + (a_n + b_n)v_n) \\
             &= (a_1 + b_1)w_1 + \cdots + (a_n + b_n)w_n \\
             &= (a_1 w_1 + \cdots + a_n w_n) + (b_1 w_1 + \cdots + b_n w_n)  \\
             &= T(u) + T(v),
  \end{align}

  and similarly,
  \begin{align}
    T(\lambda v) &= T(\lambda b_1 v_1 + \cdots + \lambda b_n v_n)  \\
                 &= \lambda b_1 w_1 + \cdots + \lambda b_n w_n  \\
                 &= \lambda (b_1 w_1 + \cdots + b_n w_n) \\
                 &= \lambda T(v).
  \end{align}

  Hence, $T$ is linear, so $T \in \linmap(V,W)$.

  Finally, to prove uniqueness, suppose \( S \in \linmap(V, W) \) is any linear map such that \( S(v_k) = w_k \) for each $k=1, \ddd, n$. Then, for any \( v = c_1 v_1 + \cdots + c_n v_n \in V \), we have
  \begin{align}
    S(v) &= c_1 S(v_1) + \cdots + c_n S(v_n) \\
         &= c_1 w_1 + \cdots + c_n w_n \\
         &= T(v),
  \end{align}

  so \( S = T \). Therefore, the unique linear map sending each $v_k$ to $w_k$ (for $k \in \{1, \ddd, n \}$) is $T$.
\end{prf}

% 3.5
\setcounter{thm}{4}
\begin{mydef} [addition and scalar multiplication on $\linmap(V,W)$]
  Suppose $S, T \in \lvw$ and $\lambda \in \mathbb{F}$. \\
  The sum $S+T \in \linmap (V,W)$ and the product $\lambda T \in \linmap (V,W)$ are defined by:
  \begin{itemize}
    \item $(S+T)(v) := Sv+Tv \quad \forall v \in V$
    \item $(\lambda T)(v) : = \lambda (Tv) \quad \forall v \in V$
  \end{itemize}
\end{mydef}

% 3.6
\subsection{Algebraic Operations on \texorpdfstring{$\linmap(V,W)$}{L(V,W)}}
\setcounter{thm}{5}
\begin{thm} [$\lvw$]
  With these operations of addition and scalar multiplication as defined above, $\lvw$ is a \vs.
\end{thm}

% 3.7
\setcounter{thm}{6}
\begin{mydef}
  If $T \in \lin{U}{V}$ and $S \in \lin{V}{W}$, then the \qt{product} $ST \in \lin{U}{W}$ is defined by
  \begin{itemize}
    \item[] $(ST)(u) := S(Tu) \quad \forall u \in U$.
  \end{itemize}

  Thus, $ST$ is just the usual composition $S \circ T$ of two functions. But with linear maps, we have distributive and associative properties (see next result), so the notation seems natural.
\end{mydef}

% 3.8
\begin{thm} [algebraic properties of products of linear maps]
  \label{thm: algebraic properties of products of linear maps}
  With these definitions, we have:
  \begin{itemize}
    \item \bfemph{associativity:} $(T_1 T_2) T_3 = T_1 (T_2 T_3)$, whenever $T_3$ maps into the domain of $T_2$ and $T_2$ maps into the domain of $T_1$.
    \item \bfemph{identity:} $T I = I T = T$ for $T \in \lvw$ (The first $I$ is the identity operator on $V$ and the second $I$ is the identity operator on $W$. We could also write $T I_V = I_W T$).
    \item \bfemph{distributive properties:} For $T, T_1, T_2 \in \lin{U}{V}$ and $S, S_1, S_2 \in \lin{V}{W}$: \\ $(S_1 + S_2)T=S_1 T + S_2 T$ and $S(T_1 + T_2)=S T_1 + S T_2$.
    \item \bfemph{non-commutative:} $ST \neq TS$ in general.
  \end{itemize}
\end{thm}

% 3.10
\setcounter{thm}{9}
\begin{thm} [linear maps take $0$ to $0$]
  \label{thm: linear maps take 0 to 0}
  $T\in \lvw \implies T(0)=0.$
\end{thm}
\begin{prf}
  $T(0) = T(0+0) = T(0) + T(0)$. Now subtract $T(0)$ on both sides.
\end{prf}